\documentclass{article}
\usepackage[UTF8]{ctex}
\usepackage{amsmath,amssymb,amsfonts}
\usepackage{enumitem}
\usepackage{geometry}
\geometry{a4paper, margin=1in}
% === 错题本排版优化（不参与题目标识） ===
\setlength{\parindent}{0pt}
\setlength{\parskip}{0.35em plus 0.10em minus 0.08em}
\setlength{\emergencystretch}{3em}
\setlist{itemsep=0.25em, topsep=0.40em, parsep=0pt}
\allowdisplaybreaks[2]
% === 错题本排版优化结束 ===
\begin{document}
\section*{福州大学2026年数学分析真题}
\section*{118．福州大学2026年数学分析试卷}

\subsection{题目1}
一．（20 分）计算极限 \(\lim _{x \rightarrow 1} \frac{e^{-\frac{(x-1)^{2}}{2}}-\cos (x-1)}{(x-1)^{2} \sqrt{2 x-x^{2}}}\) ．

\subsection{题目6}
二．（20 分）证明闭区间 \([a, b]\) 上的连续函数 \(f(x)\) 存在最大值和最小值．

\subsection{题目7}
三．（20 分）设函数 \(f(x)\) 在 \([a,+\infty)\) 上一致连续，且 \(\int_{a}^{+\infty} f(x) \mathrm{d} x\) 收敛，证明： \(\lim _{x \rightarrow+\infty} f(x)=0\) ．

\subsection{题目2}
四．（20分）求级数 \(\sum_{n=1}^{\infty} a_{n} x^{n}\) 的收敛域与和函数，其中 \(a_{n}=1-\frac{1}{2}+\frac{1}{3}-\cdots+(-1)^{n-1} \frac{1}{n}\) 。

\subsection{题目3}
五．（20分）计算由曲线 \(y^{3}=x, y^{3}=2 x\) 以及 \(x y=1, x y=2\) 所围区域在第一象限部分的面积．

\subsection{题目4}
六．（25 分）计算单位时间内不可压缩流体（密度为 1 ）以速度 \(v=x^{2} i+y^{2} j+z^{2} k\) 流过曲面 \(\sum\) 的流量，其中\\
\(\sum: x^{2}+y^{2}=z^{2}, 0 \leq z \leq 1\), 指向下侧。

\subsection{题目5}
七．（25 分）讨论函数 \(I(\alpha)=\int_{0}^{+\infty} \frac{\cos x^{2}}{x^{\alpha}}\) 在 \(\alpha \in(-1,1)\) 的连续性．
\end{document}
